% Source: User reference p. 105 reroot framework; distance-sum formula derived explicitly.
\subsection{换根 DP：所有点的距离和}
\label{dp:reroot-dp}
\AlgorithmInfo
    {求树上每个顶点到所有顶点的加权距离和。}
    {时间、空间 $O(n)$。}
    {非空连通树、1 下标；边权非负，所有相关距离和在 64 位范围内。}
\subsubsection*{原理}

第一次递归 DFS 求子树大小 $sz_u$ 和 $down_u$，后者为 u 到其子树各点的距离和：
\[
 down_u=\sum_{v\text{ 为孩子}}(down_v+sz_v w_{uv}).
\]
叶子的 sz 为 1、down 为 0；先处理孩子再合并，初始根的答案就是它的 down。根沿边从 u 移到孩子 v 时，v 子树中的 $sz_v$ 个点都近 $w$，其余 $n-sz_v$ 个点都远 $w$，所以
\[
 ans_v=ans_u+(n-2sz_v)w_{uv}.
\]
第二次 DFS 将根外贡献向下传递，即可在线性时间内得到所有根的答案。其他换根问题也要明确如何移除该孩子的贡献，再合并父侧贡献。

\subsubsection*{使用教程}
\InlineCode{distance_sums(e,root)} 接收加权双向邻接表，返回按 1 下标编号的距离和数组；预处理根默认 1，改变它不改变最终答案。每条无向边需要在两个端点各存一次；无权树将权值填为 1。
\begin{lstlisting}[style=mtycpp,numbers=none]
vector<vector<pair<int,long long>>> e(5);
auto add=[&](int u,int v,long long w){
	e[u].push_back({v,w});e[v].push_back({u,w});
};
add(1,2,3);add(2,3,2);add(2,4,1);
auto ans=distance_sums(e);
long long sum=ans[3];
\end{lstlisting}
\needspace{7\baselineskip}
\subsubsection*{模板代码}
\CodeFile{code/dp/reroot_dp.cpp}
